\documentclass[12pt,a4paper]{article}

\usepackage[utf8]{inputenc}
\usepackage[T1]{fontenc}
\usepackage[spanish,es-tabla]{babel}
\usepackage{amsmath,amssymb}
\usepackage{geometry}
\usepackage{siunitx}
\usepackage{circuitikz}

\geometry{margin=2.5cm}

\title{Problema 4-1: Función de transferencia de un circuito RC}
\author{}
\date{}

\begin{document}

\maketitle

\section*{Planteamiento}

El circuito está formado por una resistencia $R_1$ en serie con una rama constituida
por un capacitor $C$ y una resistencia $R_2$ conectados en serie. La tensión de entrada
es $E_1(s)$ y la tensión de salida $E_2(s)$ se mide sobre la rama $C$--$R_2$.

\begin{center}
\begin{circuitikz}[american voltages]
    \draw
    (0,0) to[voltage source,l=$E_1(s)$] (0,3)
    (0,3) to[R,l=$R_1$] (3,3)
    (3,3) to[C,l=$C$] (3,1.5)
    (3,1.5) to[R,l=$R_2$] (3,0)
    (3,0) -- (0,0);
    \draw (3.4,0) to[open,v^=$E_2(s)$] (3.4,3);
\end{circuitikz}
\end{center}

\section*{Desarrollo}

La impedancia del capacitor en el dominio de Laplace es

\begin{equation}
    Z_C(s)=\frac{1}{sC}.
\end{equation}

Por tanto, la impedancia de la rama donde se mide la salida es

\begin{equation}
    Z_{CR_2}(s)=R_2+\frac{1}{sC}.
\end{equation}

Aplicando la regla del divisor de tensión:

\begin{equation}
    \frac{E_2(s)}{E_1(s)}
    =\frac{Z_{CR_2}(s)}{R_1+Z_{CR_2}(s)}.
\end{equation}

Sustituyendo la impedancia de la rama RC:

\begin{align}
    H(s)&=\frac{E_2(s)}{E_1(s)}\\[4pt]
    &=\frac{R_2+\dfrac{1}{sC}}
            {R_1+R_2+\dfrac{1}{sC}}.
\end{align}

Multiplicando numerador y denominador por $sC$, se obtiene la función de transferencia:

\begin{equation}
    \boxed{
    H(s)=\frac{E_2(s)}{E_1(s)}
    =\frac{1+sCR_2}{1+sC(R_1+R_2)}
    }.
\end{equation}

\section*{Polo y cero}

El cero se obtiene anulando el numerador:

\begin{equation}
    s_z=-\frac{1}{CR_2}.
\end{equation}

El polo se obtiene anulando el denominador:

\begin{equation}
    s_p=-\frac{1}{C(R_1+R_2)}.
\end{equation}

Una forma alternativa de escribir la función de transferencia es

\begin{equation}
    H(s)=
    \frac{R_2}{R_1+R_2}
    \frac{1+sCR_2}
         {1+sC\dfrac{R_1R_2}{R_1+R_2}}.
\end{equation}

\section*{Comportamiento en frecuencia}

En corriente continua ($s=0$), el capacitor se comporta como un circuito abierto:

\begin{equation}
    H(0)=\frac{R_2}{R_1+R_2}.
\end{equation}

Para frecuencias muy altas ($s\rightarrow\infty$), el capacitor se comporta como un
cortocircuito:

\begin{equation}
    \lim_{s\rightarrow\infty}H(s)=1.
\end{equation}

Por ello, el circuito no es un pasa-altos ideal. Es una red de adelanto o una red de
transici\'on, cuya ganancia aumenta desde $R_2/(R_1+R_2)$ en baja frecuencia hasta $1$
en alta frecuencia.

\end{document}
